% ============================================================================
%   CHAPTER 6 — CONCLUSION AND FUTURE WORK
%   Aivancity: 3-5 pages
%   REVISED to integrate the decoding-artefact reinterpretation, the
%   re-thresholding remedy, the collapse/corruption taxonomy, and the
%   TREC external replication.
% ============================================================================

\chapter{Conclusion and Future Work}
\label{chap:conclusion}

% ----------------------------------------------------------------------------
\section{Summary of Findings}
\label{sec:conc-summary}

This thesis reports a systematic calibration audit of twelve open-source 
large language model configurations applied to patient--clinical trial 
matching. The study was structured around four pre-registered hypotheses, 
tested with fixed rejection thresholds and Bonferroni-corrected 
statistical procedures, and complemented by twelve original findings that 
emerged from the deeper analysis. Table~\ref{tab:res-summary-hypotheses} 
in Chapter~\ref{chap:results} summarizes the pre-registered outcomes 
and Table~\ref{tab:res-summary-findings} lists the twelve findings.

The pre-registered results support the central hypothesis of the study 
and challenge one of its ancillary claims. Hypothesis~H1 is confirmed 
across all twelve configurations: every LLM in the experimental matrix 
exhibits an Expected Calibration Error above the well-calibrated 
threshold of $0.05$~\cite{vancalster2019calibration}, with values 
ranging from $0.124$ to $0.602$. Hypothesis~H2 is confirmed on both 
model families tested (Mistral-7B and Llama-3.1-8B) and both prompt 
variants: RLHF-aligned Instruct versions exhibit significantly worse 
calibration than their Base counterparts, all four permutation tests 
rejecting the null hypothesis at Bonferroni-corrected significance 
($p < 0.003$). This is the central pre-registered contribution of the 
thesis, and it extends to the clinical trial matching domain a finding 
first documented by Kadavath et al.~\cite{kadavath2022know} in 
general-purpose question answering. Hypothesis~H3 is not supported: no 
configuration reaches the C-index threshold of $0.70$, though Meditron 
approaches it at $0.654$. Hypothesis~H4 is partially supported, with 
the selective-classification gain reaching the $+0.02$ threshold only 
for Mistral configurations.

Beyond the pre-registered hypotheses, the exploratory findings fall into 
two waves. The first wave documents the failure modes themselves. The 
enriched-prompt paradox (F3) shows that expert prompt engineering 
degrades calibration in five of six models, contradicting the standard 
recommendation to invest in prompt design as a route to better model 
behavior. Mode collapse (F7 and F8) affects five of the twelve 
configurations, which concentrate more than $97\%$ of their predictions 
on a single label and do not recover the diversity of their predicted 
labels under constrained decoding. Complete failure on truly 
discriminative cases (F9) shows that both medical LLMs tested (Meditron 
and BioMistral) achieve accuracy at or below $0.4\%$ on the $238$ pairs 
requiring a genuine clinical decision, despite their apparent 
superiority on global metrics.

The second wave, developed through a mechanistic analysis of the 
underlying log-probabilities, qualifies and explains these failures. The 
mode collapse of four of the five affected configurations is not an 
absence of clinical knowledge but a property of the $\arg\max$ decision 
rule: the discriminative signal is present in the probabilities --- 
Llama-Instruct~A obtains the highest ordinal separation of the entire 
study while predicting a single label $97.8\%$ of the time --- but its 
amplitude is too small to overcome a response prior installed by 
post-training. A single scalar, the ratio of signal amplitude to prior 
amplitude, orders the twelve configurations by their collapse behaviour 
(F11) and explains, as a mathematical consequence, why monotone 
recalibration cannot repair collapse while a non-monotone re-thresholding 
rule can: the latter restores between $+0.62$ and $+0.79$ macro-F1 on the 
collapsed configurations that carry a live signal, with no gain on the 
non-collapsed ones (F10). The same analysis separates recoverable 
collapse (BioMistral) from an irrecoverable corruption of the signal 
(Meditron, whose discriminative AUC lies below chance), a distinction 
that aggregate accuracy had merged under Finding~F9. Finally, both the 
collapse signature and its signal/prior mechanism replicate on the 
independent TREC Clinical Trials corpus under an identical prompt and 
scoring pipeline (F12), confirming that these are properties of the 
models rather than artefacts of the original benchmark.

% ----------------------------------------------------------------------------
\section{Answer to the Research Question}
\label{sec:conc-answer}

The research question posed in Chapter~\ref{chap:introduction} asked 
whether the eligibility scores produced by LLM-based patient--trial 
matching systems are probabilistically reliable, and whether RLHF 
post-training affects this reliability. The answer, informed by the 
twelve configurations audited in this study, is negative on both 
dimensions as the systems stand --- but the mechanistic analysis 
qualifies the nature of the failure.

The scores are not reliable as deployed: the systematic miscalibration 
documented by Hypothesis~H1 disqualifies every configuration from direct 
deployment as an autonomous eligibility classifier, and the mode 
collapse of Findings F7 and F8 disqualifies five of the twelve 
configurations from any useful classification role under the default 
$\arg\max$ rule. Only Mistral-7B-Instruct with the minimal Prompt~A 
reaches clinically acceptable performance on the truly discriminative 
subset without further intervention.

This unreliability, however, is not uniformly a matter of missing 
knowledge. For four of the five collapsed configurations, the 
discriminative signal is present but masked by the decision rule; the 
unreliability is a decoding failure that a threshold rule recovers, not 
an absence of clinical competence. The distinction matters for 
deployment: a model in the expressive-collapse regime is unusable 
off-the-shelf yet salvageable with a lightweight, weight-free 
intervention, whereas a model in the corruption regime (Meditron) ranks 
cases below chance and is salvageable by no post-hoc rule. The answer to 
the research question is therefore layered: the scores are not reliable, 
the RLHF-aligned scores are the least reliable of all, but the latent 
discriminative information is, for several models, both present and 
recoverable.

RLHF post-training significantly worsens calibration on both model 
families tested, converting the underconfidence of Mistral-Base into 
the overconfidence of Mistral-Instruct, and inducing severe 
overconfidence in Llama-Instruct. The effect is quantitatively 
amplified relative to the general-QA setting studied by Kadavath 
et al., with $\Delta\ece$ values reaching $+0.451$ in the worst case 
--- approximately three times the magnitude previously reported. The 
mechanistic analysis refines the account of \emph{how} RLHF does this: 
it amplifies both the discriminative signal and a response prior, and 
degrades calibration --- or, in the extreme, induces collapse --- 
precisely when it raises the prior faster than the signal. In the Llama 
family this ratio falls below one and the model collapses; in the 
Mistral family it rises and discrimination is preserved. RLHF is thus 
not uniformly harmful to reliability; it is harmful when it makes the 
response prior dominate the signal the model still holds.

Post-hoc isotonic recalibration (F6) provides a partial remedy for the 
models that classify: it brings eleven of the twelve configurations 
below the calibration threshold, robustly across five-fold 
cross-validation. It cannot, by construction, help the collapsed models, 
because it is monotone and preserves their frozen $\arg\max$. The 
complementary remedy for collapse is re-thresholding (F10), which is 
mathematically capable of shifting the decision and empirically restores 
substantial macro-F1 --- but only for the configurations whose signal 
survives, identified by an above-chance AUC. The practical conclusion is 
that a pre-deployment audit should combine both remedies with a 
diagnostic that decides which applies: recalibration for the diverse 
predictors, re-thresholding for the recoverable collapses, and rejection 
for the corrupted models.

% ----------------------------------------------------------------------------
\section{Contributions}
\label{sec:conc-contributions}

The scientific contributions of this thesis fall into four groups.

First, the study provides the first systematic calibration audit of 
open-source LLMs applied to clinical trial matching. The audit follows 
the reporting recommendations of TRIPOD+AI~\cite{collins2024tripod} and 
covers twelve configurations organized in a two-by-two experimental 
matrix (generalist versus medical, Base versus Instruct). Every 
prediction and every numerical result is archived in a public 
repository to support independent replication (Appendix~\ref{app:code}).

Second, the study extends the RLHF-degrades-calibration finding of 
Kadavath et al.~\cite{kadavath2022know} to the clinical domain, with a 
pre-registered permutation test replicated on two independent model 
families under two independent prompt designs. To our knowledge, this 
is the first replication of this effect outside general-purpose 
question answering, and the first quantification of its amplification 
in a high-stakes clinical task.

Third, the study documents the failure modes of open-source medical 
LLMs that had not been previously reported in a systematic form: the 
enriched-prompt paradox (F3), mode collapse and its persistence under 
constrained decoding (F7, F8), and complete failure on discriminative 
cases (F9). These findings jointly motivate the composite utility score 
introduced in Section~\ref{subsec:met-composite}, which combines 
calibration, discrimination, and prediction diversity into a single 
rank-ordering that resists majority-class shortcuts.

Fourth --- and this is the contribution that emerged from the deeper 
analysis --- the study reinterprets mode collapse as a decoding 
artefact rather than a knowledge loss, and turns this reinterpretation 
into three usable results. It identifies a scalar mechanism, the 
signal-to-prior ratio, that predicts collapse across configurations and 
explains why recalibration fails where re-thresholding succeeds. It 
proposes a two-dimensional diagnostic (signal/prior ratio crossed with 
discriminative AUC) that separates recoverable collapse from 
irrecoverable corruption, and shows that the re-thresholding F1 gain is 
a misleading recovery indicator unless paired with an above-chance AUC. 
And it validates both the collapse signature and its mechanism on the 
independent TREC Clinical Trials corpus. Together these turn a 
descriptive failure mode into a diagnosable, and often correctable, 
property of the decision rule.

The methodological byproducts of the study --- the pre-registered 
protocol, the double-prompt design, the log-probability scoring 
pipeline, the composite utility score, and the signal/prior diagnostic 
--- are intended as reusable elements for future calibration audits of 
clinical LLMs.

% ----------------------------------------------------------------------------
\section{Limitations}
\label{sec:conc-limitations}

Three categories of limitation qualify the conclusions of this work: 
the reliance on a single primary benchmark, the constrained model 
scale, and the finite prompt space explored.

\subsection{External Validation}
\label{subsec:conc-limitations-external}

The primary analyses of this study rest on a single benchmark, 
TrialGPT-Criterion-Annotations~\cite{jin2024trialgpt}, chosen because 
it is the only public corpus offering per-criterion multi-level expert 
annotations with an explicit inclusion/exclusion distinction --- 
properties that are essential for the fine-grained calibration analysis 
and the criterion-level signal/prior mechanism. To test whether the 
central mechanistic findings depend on this benchmark, we conducted an 
external replication on the TREC Clinical Trials 2021--2022 corpus 
(Section~\ref{sec:res-external}), using an identical prompt and scoring 
pipeline and changing only the corpus. The collapse signature and the 
signal/prior ratio ordering reproduced, and the recoverable-collapse 
versus corruption distinction persisted at the trial level (F12). This 
replication substantially strengthens the external validity of the 
mechanistic contributions.

Two limitations of this external replication should nonetheless be 
stated plainly. First, TREC provides relevance judgements at the trial 
level and on a three-grade scale, not the six-level per-criterion 
annotations of TrialGPT-Criterion; the trial-level recoverability we 
report on TREC therefore required an aggregation step and is a 
convergent, partial validation rather than an exact replication of the 
criterion-level recovery magnitude. The pre-registered calibration 
hypotheses H1--H4, which depend on the six-level scale, were not 
re-tested on TREC. Second, external validation on real, prospectively 
collected clinical data --- as opposed to the synthetic or 
retrospective notes of both benchmarks --- would require ethical 
approval, hospital partnership, and integration into a live recruitment 
workflow, and remains a research project in its own right. The 
mechanistic replication reported here closes the most immediate gap 
(specificity to a single benchmark) but does not substitute for 
prospective clinical validation.

\subsection{Model Scale}
\label{subsec:conc-limitations-scale}

The models evaluated in this study are limited to 7--9 billion 
parameters, which is the size range that fits into a free-tier T4 GPU 
with 4-bit quantization. Larger medical LLMs, such as Meditron-70B or 
Med42-v2-70B, may exhibit qualitatively different calibration and 
failure-mode patterns. In particular, larger models may carry a larger 
signal relative to their response prior, which under the mechanism of 
Section~\ref{subsec:res-collapse-mechanism} would raise the 
signal/prior ratio and attenuate collapse; whether this occurs is an 
empirical question the present budget could not address. The 
extrapolation of the results to larger scales is a natural direction 
for future work, and the methodology --- including the signal/prior 
diagnostic --- is directly transferable to 70B-scale configurations 
provided that the computational budget is available.

\subsection{Prompt Space and Prompt-Driven Effects}
\label{subsec:conc-limitations-prompts}

Only two prompt variants were evaluated in the present study: a 
minimal Prompt~A and an enriched Prompt~B, both reproduced verbatim 
in Appendix~\ref{app:prompts}. The space of possible prompts is much 
larger, encompassing chain-of-thought designs, role-based framings, 
and few-shot examples selected through active learning. The 
enriched-prompt paradox (F3) documented here is thus specific to the 
two prompts tested; it may or may not generalize to other prompt 
families. A systematic study of the calibration effect of prompt 
design is a promising direction for follow-up work.

The choice of prompt is not neutral, and it is legitimate to ask 
whether the failure modes documented here are properties of the 
models or artifacts of the prompt design. Four observations argue 
against a purely prompt-driven interpretation. First, Hypothesis~H2 
is tested through paired comparisons within the same prompt (Base 
versus Instruct on Prompt~A, Base versus Instruct on Prompt~B), 
which cancels out any systematic effect of the prompt on the ECE 
difference. Second, mode collapse affects only five of the twelve 
configurations, whereas a prompt-driven artifact would be expected to 
affect all models exposed to that prompt in a similar way; the 
observed heterogeneity across models sharing the same prompt is 
incompatible with a single prompt-level cause. Third, the collapse is 
governed by an internal quantity --- the signal-to-prior ratio 
(Section~\ref{subsec:res-collapse-mechanism}) --- that varies across 
models under the same prompt and predicts the collapse independently of 
prompt content. Fourth, the collapse signature and its mechanism 
replicate on an entirely different corpus under the same prompt (F12), 
which a prompt artefact specific to TrialGPT-Criterion could not do. 
That said, the specific proportion of collapsed configurations, the 
magnitude of the enriched-prompt paradox (F3), and the direction of the 
inclusion/exclusion asymmetry (F4) may depend on the particular prompts 
tested here. Their exact quantitative values should be interpreted as 
characteristic of the present experimental setup rather than as 
universal constants of the underlying models.

% ----------------------------------------------------------------------------
\section{Future Work}
\label{sec:conc-future}

Four research directions extend naturally from the present study, in 
order of feasibility.

\subsection{Broader and Prospective External Validation}
\label{subsec:conc-external}

The external replication on TREC (Section~\ref{sec:res-external}) 
establishes that the collapse mechanism is not specific to 
TrialGPT-Criterion. Three extensions would broaden this validation 
further. First, re-testing the pre-registered calibration hypotheses 
H1--H4 on a corpus with per-criterion multi-level annotations, were one 
to become available, would extend the confirmatory results as well as 
the mechanistic ones. Second, MIMIC-derived cohorts and n2c2 datasets, 
which contain more heterogeneous and noisier patient descriptions, would 
test whether the signal/prior mechanism survives realistic clinical 
documentation styles. Third, and most demanding, prospective validation 
on live recruitment workflows would close the gap between benchmark 
performance and clinical utility.

\subsection{Robust Re-thresholding and Prior Correction}
\label{subsec:conc-rethreshold}

The re-thresholding remedy of Section~\ref{sec:res-rethreshold} 
recovers substantial macro-F1 but leaves two questions open. The 
threshold itself is tuned per configuration and, for one collapsed 
model, proved unstable across folds; a prior-correction rule that 
subtracts the estimated response prior from the log-probabilities before 
the $\arg\max$ --- a multi-class generalization of the binary threshold 
--- would remove the free parameter and may transfer more robustly 
across datasets. Characterizing the conditions under which such a 
correction generalizes, and comparing it to the per-criterion threshold 
used here, is a concrete and immediately feasible extension.

\subsection{Multi-Agent Adversarial Debate for Calibration}
\label{subsec:conc-debate}

The failure of individual models on discriminative cases motivates 
investigating whether ensembles of models can outperform any single 
one. Multi-agent debate architectures, in which several LLMs exchange 
arguments before converging on a joint prediction 
\cite{ouyang2022training}, are a particularly interesting candidate: 
models with distinct RLHF biases and distinct evasion channels may 
mutually neutralize each other's overconfidence through argumentative 
confrontation. The 6-model consensus analysis reported in 
Section~\ref{subsec:res-consensus}, which shows a $75.5\%$ accuracy on 
the $163$ pairs of full agreement, suggests that this direction is 
promising. We note, however, that the re-thresholding results temper the 
motivation: several models that appear useless under consensus voting 
in fact carry a recoverable individual signal, so a debate among 
re-thresholded models may be a more informative baseline than a debate 
among their raw $\arg\max$ outputs.

\subsection{Investigating the Origin of the Response Prior}
\label{subsec:conc-collapse-origin}

The mechanistic analysis of this thesis identifies the \emph{proximal} 
cause of mode collapse --- a response prior that dominates the 
discriminative signal --- but leaves open its \emph{distal} cause. Why 
does RLHF install a prior that grows faster than the signal in the 
Llama family but not in the Mistral family, and why does continued 
medical pretraining corrupt the signal outright in Meditron? Answering 
this would require access to the training corpora and reward models of 
the affected configurations, and interpretability analyses of how the 
response prior is represented internally. Because the signal/prior ratio 
gives a concrete, measurable target, such an investigation could be 
framed as identifying which training interventions raise the prior 
relative to the signal --- a sharper question than the original, purely 
descriptive one of what causes collapse.

\subsection{Extension to Larger and Multilingual Models}
\label{subsec:conc-scale-future}

Extending the analysis to 70-billion-parameter models is a 
straightforward but resource-intensive direction. Meditron-70B and 
Med42-v2-70B are the natural targets, and the signal/prior diagnostic 
provides a direct prediction to test: if larger models resist collapse, 
their signal/prior ratio should exceed that of their 7B counterparts. 
Multilingual evaluation is a second direction: models such as 
MMed-Llama-3 and Apollo have been developed to support clinical text in 
languages other than English, and the calibration properties of these 
multilingual models have not been systematically audited on any 
language.

% ----------------------------------------------------------------------------
\section{Closing Remarks}
\label{sec:conc-closing}

The systematic use of large language models in clinical prediction 
workflows is likely to expand rapidly in the coming years, driven by 
the maturation of open-source alternatives, the pressure to accelerate 
clinical trial recruitment, and the availability of increasingly 
capable models. This trajectory makes the auditing of these models a 
practical necessity rather than an academic exercise. The present 
thesis contributes to this auditing effort in four specific ways: it 
proposes a reusable methodology, it documents previously unreported 
failure modes in open-source medical LLMs, it reinterprets the most 
severe of these failures as a correctable property of the decision 
rule, and it archives its data and code for independent verification.

The broader lesson of this study is that global performance metrics 
--- accuracy, F1, and even ECE taken in isolation --- are insufficient 
to characterize the clinical usability of a language model, and that 
the predicted-label distribution can itself be misleading: a model that 
predicts a single label may nonetheless carry a strong, recoverable 
discriminative signal underneath it. A responsible evaluation protocol 
must therefore combine calibration with diversity, must stratify 
performance by case difficulty, must remain vigilant against mode 
collapse, and must distinguish a collapse of the decision rule from a 
corruption of the signal --- because the first is correctable and the 
second is not.

None of the twelve open-source configurations audited here provides a 
fully satisfactory basis for autonomous clinical trial matching as 
deployed off the shelf. The findings do, however, identify 
Mistral-Instruct with a minimal prompt as the most promising candidate 
for direct use, identify post-hoc isotonic recalibration as a practical 
remedy for the miscalibration of the models that classify, and identify 
re-thresholding as a complementary remedy that restores substantial 
discriminative performance to the models whose collapse masks a live 
signal. Whether these elements are jointly sufficient for a deployment 
scenario --- with mandatory human oversight, ongoing calibration 
monitoring, a signal/prior screen to route each model to the 
appropriate remedy, and protocols for handling low-confidence cases --- 
is a question that the present thesis leaves to prospective clinical 
validation. Answering it is the natural next step of this line of 
research.